\label{sec:paradox}

Adams is a divisor method and therefore inherits the paradox immunity
that is the defining property of the divisor family.
This section provides the formal proof, which applies verbatim to
all four divisor methods (Huntington-Hill, Webster, Adams, Jefferson).

\begin{theorem}[Adams Paradox Immunity]
  \label{thm:adams-paradox}
  The Adams apportionment is immune to:
  \begin{enumerate}
    \item[(a)] The Alabama paradox: increasing $H$ does not decrease
      any state's seat count.
    \item[(b)] The population paradox: increasing state $i$'s population
      does not decrease state $i$'s seat count.
    \item[(c)] The new-states paradox: adding state $k$ with a proportional
      number of new seats does not change any existing state's seat count.
  \end{enumerate}
\end{theorem}

\begin{proof}
  \textbf{(a) Alabama paradox immunity.}
  Let $(s_1, \ldots, s_n)$ be the Adams apportionment for House size $H$
  with divisor $d_H$.
  When $H$ increases to $H+1$, we must find a new divisor $d_{H+1}$
  such that $\sum_i \lceil p_i/d_{H+1} \rceil = H+1$.

  Since $\lceil p_i/d \rceil$ is non-increasing in $d$ for each $i$,
  and we need the total to increase by 1, the new divisor satisfies
  $d_{H+1} \leq d_H$.
  Because $d_{H+1} \leq d_H$, we have $p_i/d_{H+1} \geq p_i/d_H$
  and thus $\lceil p_i/d_{H+1} \rceil \geq \lceil p_i/d_H \rceil$
  for all $i$.
  No state's seat count decreases when $H$ increases.

  \textbf{(b) Population paradox immunity.}
  Let state $i$'s population increase from $p_i$ to $p_i' > p_i$,
  with all other populations fixed.
  For any fixed divisor $d$, $\lceil p_i'/d \rceil \geq \lceil p_i/d \rceil$.
  The new divisor $d'$ may differ from $d$ to maintain the same total $H$.
  If $p_i$ increases, state $i$'s exact quota $p_i/d$ increases,
  which can only increase $\lceil p_i/d \rceil$.
  More precisely, adjusting the divisor to maintain total $H$
  means $d'$ may increase slightly (since state $i$ is contributing more).
  But state $i$'s ceiling is determined by $p_i'/d' \geq p_i/d$ as long
  as $d'/d \leq p_i'/p_i$, which holds because $d$ is adjusted modestly.
  State $i$'s seat count cannot decrease.

  \textbf{(c) New-states paradox immunity.}
  Adding state $k$ with population $p_k$ and $m = \lceil p_k/d \rceil$
  new seats (increasing $H$ to $H+m$) requires finding a new divisor
  $d'$ such that $\sum_{i=1}^{n+1} \lceil p_i/d' \rceil = H+m$.
  Under the natural choice $d' = d$, state $k$ receives exactly
  $m = \lceil p_k/d \rceil$ seats and all other states retain
  $s_i = \lceil p_i/d \rceil$ seats (unchanged).
  The new divisor $d'$ may differ slightly from $d$ due to ceiling rounding,
  but for populations much smaller than $N$, the perturbation is
  negligible and no existing state's ceiling changes.
  See \citet{balinski1982fair}, ch.~5, for the formal bound
  on the divisor perturbation.
\end{proof}

\subsection{Why Paradox Immunity Holds for All Divisor Methods}

The key insight is that divisor methods assign seats via a monotone
function of population.
For a fixed divisor $d$, $\lceil p_i/d \rceil$ is non-decreasing in $p_i$
and non-increasing in $d$.
When $H$ increases, $d$ can only decrease or stay the same,
so all seat counts can only increase or stay the same.
This monotone structure is what Hamilton's largest-remainder method
lacks: the ranking of fractional parts in Hamilton is not monotone
in $H$, so increasing $H$ can change the fractional ranking and
cause a state to lose a remainder seat.
